\documentclass[../../main.tex]{subfiles}
\begin{document}

\begin{center}
  {\fontsize{48pt}{58pt}\selectfont \textbf{1985}\normalsize\textbf{年度}}
\end{center}

\begin{table}[htb]
  \centering
  \caption{出題テーマと難易度}
  \begin{tabular}{|c|p{9cm}|c|}
    \hline
    \textbf{問題番号} & \textbf{テーマ} & \textbf{難易度} \\
    \hline
    第1問 & 正弦曲線で定まる領域の面積差 & \\
    \hline
    第2問 & 円周上の二点と射影間の距離 & \\
    \hline
    第3問 & 平面対称で得る四面体の体積 & \\
    \hline
    第4問 & 行列の累乗と線分長の最小化 & \\
    \hline
    第5問 & 独立な確率変数と分布の決定 & \\
    \hline
    第6問 & 球が通過する立体の断面積と体積 & \\
    \hline
  \end{tabular}
\end{table}

\begin{table}[htb]
  \centering
  \caption{解答}
  \renewcommand{\arraystretch}{1.4}
  \begin{tabular}{|c|c|p{9cm}|}
    \hline
    \textbf{問題番号} & \textbf{小問} & \textbf{答え} \\
    \hline
    第1問 & -- & $a=\dfrac{2\sqrt{3}}{3}$のとき，最大値は$\dfrac{\pi}{3}$ \\
    \hline
    第2問 & -- & $l=0$なら$3$通り，$0<l<\dfrac{9}{8}$なら$6$通り，$l=\dfrac{9}{8}$なら$4$通り，$\dfrac{9}{8}<l<2$なら$2$通り，$l=2$なら$1$通り，$l>2$なら$0$通り \\
    \hline
    第3問 & -- & $V=\dfrac{4t^4}{3(1+t^2)^2}$ \\
    \hline
    第4問 & (1) & \\
    \cline{2-3}
    & (2) & \\
    \hline
    第5問 & -- & $p_n=\left(\dfrac{1}{2}\right)^{n+1}$，$\displaystyle\sum_{n=0}^{\infty}np_n=1$ \\
    \hline
    第6問 & (1) & 面積は$21+6\sqrt{3}+\pi$ \\
    \cline{2-3}
    & (2) & 体積は$4\left(\dfrac{10}{3}+\sqrt{3}\right)\pi-4\left(1+\dfrac{2\sqrt{3}}{3}\right)$ \\
    \hline
  \end{tabular}
\end{table}

\paragraph{試験形態}
6問．試験時間は要確認．

\end{document}
